\chapter{Preference Elicitation in the Era of Large Language Models: An Updated Review}
\label{ch:llm_lit}

\section{Introduction and Scope}

The preceding chapters of this thesis reviewed the Conversational Recommender Systems (CRS) literature as it stood at the time the bot-play elicitation framework was designed, and Chapter~\ref{ch:main} presented a reinforcement learning approach in which a Questioner agent learns, through self-play against a simulated user, which questions yield the greatest information gain for a downstream recommendation model. That work was subsequently published at IJCNN 2024 \citep{makarova2024learning}. Since then, the field has undergone what is arguably its most significant transformation to date: the arrival of instruction-tuned Large Language Models (LLMs) has changed the assumptions underlying every component of a CRS, from natural language understanding and generation, through user simulation, to the recommendation logic itself.

This chapter updates the literature review in light of these developments, with three objectives. First, it consolidates the lineage of reinforcement-learning-based preference elicitation that this thesis descends from, paying particular attention to the structure of the \textit{action spaces} these systems employ, since this is where the present work departs from prior art. Second, it reviews the rapid emergence of LLM-based conversational recommendation between 2023 and 2026, and examines the evidence for what remains the central motivating claim of this thesis: that the ability to \textit{strategically} elicit user preferences — deciding what to ask, and when, so as to maximise recommendation quality — is not solved by language modelling alone. Third, it re-derives the research gap addressed by this thesis with the precision that the matured field now demands, distinguishing carefully between mechanisms that constitute genuine prior art and the specific combination of techniques that remains unexplored.

The review in this chapter was conducted by first screening titles and abstracts of candidate works published between 2019 and 2026, with emphasis on 2023--2026, followed by full-text reading of the works closest to the thesis topic. Claims about the internal design of the systems discussed below (in particular, their action space structure) were verified against the primary sources rather than secondary descriptions, since, as will become apparent, the precise formulation of the decision problem is the crux of the comparison. The chapter is organised as follows. Section~\ref{sec:problem} restates strategic preference elicitation as a sequential decision problem, providing the vocabulary used throughout. Section~\ref{sec:rl_lineage} traces the discrete-action reinforcement learning lineage from its foundations to its most recent hierarchical refinements. Section~\ref{sec:continuous} reviews continuous-action reinforcement learning and large-scale action grounding, the mechanism family this thesis builds upon. Section~\ref{sec:llm_crs} surveys LLM-based conversational recommendation, and Section~\ref{sec:elicitation_llm} the parallel revival of preference elicitation research around LLMs. Section~\ref{sec:eval_norms} examines how evaluation practice has changed and identifies a methodological gap that this thesis addresses directly. Section~\ref{sec:gap} synthesises the updated research gap and states the research questions for the remainder of the thesis.

\section{Strategic Preference Elicitation as a Sequential Decision Problem}
\label{sec:problem}

Throughout this chapter it is useful to have a single formal vocabulary against which heterogeneous systems can be compared. Following the framing introduced in Chapters~\ref{ch:rl_crs} and~\ref{ch:main}, strategic preference elicitation can be cast as a finite-horizon sequential decision problem. At each conversational turn $t$, the system holds an internal representation $s_t$ of what has been learned about the user so far — variously implemented as a belief state, a dialogue state, a set of confirmed attribute values, or a dense embedding of the conversation history. The system selects an elicitation action $a_t$ from some action space $\mathcal{A}$ — a question to ask, in whatever form the system's interface supports — and the user's response is incorporated into the next state $s_{t+1}$. A recommendation model $\phi$ maps the state to a ranking over the item catalogue, and the quality of that ranking, measured against held-out ground truth, defines the utility of the elicitation so far. A questioning \textit{strategy}, or policy $\pi(a_t \mid s_t)$, is judged by how rapidly it increases recommendation utility per turn, under the practical constraint that users tolerate only a small number of questions before disengaging \citep{sun2018conversational, lei2020estimation}.

Three design choices differentiate essentially every system reviewed in this chapter. The first is the structure of $\mathcal{A}$: whether the policy chooses from a fixed enumeration of catalogue attributes, from a small set of abstract dialogue strategies, from the open space of natural language, or — as in this thesis — from a continuous semantic embedding space grounded to discrete entities. The second is whether $\pi$ is \textit{learned} by optimising a multi-turn objective, computed myopically from a one-step criterion such as entropy or expected information gain, or not optimised at all (as in prompt-only LLM systems). The third is how the chosen action is \textit{realised} as language: as a template, as a symbolic exchange with no language at all, or by a generative language model. Holding this vocabulary fixed makes it possible to see past superficial differences — between, say, a 2018 dialogue agent and a 2025 LLM pipeline — to the underlying decision problem, which has remained remarkably stable even as the surrounding machinery has been transformed.

\section{Reinforcement Learning for Preference Elicitation: The Discrete-Action Lineage}
\label{sec:rl_lineage}

\subsection{Foundational Systems: CRM and EAR}

The idea that a recommender dialogue agent should \textit{learn} its questioning behaviour, rather than follow handcrafted rules, predates the LLM era and forms a well-defined research lineage. The Conversational Recommender Model (CRM) of \citet{sun2018conversational} was among the first to integrate a neural belief tracker with a policy network trained by policy gradient: the belief tracker maintains a distribution over the values of a fixed set of item facets, and at each turn the policy decides either to request the value of one of those facets or to display a recommendation computed by a factorisation machine conditioned on the facets confirmed so far. CRM operates in a single-round setting — the dialogue ends at the first recommendation, successful or not — and was trained against simulated dialogues synthesised from templated utterances crowd-sourced at scale.

The Estimation--Action--Reflection (EAR) framework \citep{lei2020estimation} deepened the interaction between the conversational and recommendation components through its eponymous three-stage decomposition. In the \textit{estimation} stage, a factorisation machine is trained offline to score both items and attributes, conditioned on the user and the attributes confirmed during conversation; in the \textit{action} stage, a policy network receives a state vector summarising the conversation — including the entropy of the candidate item distribution and the recommender's confidence — and decides whether to ask about a further attribute or to recommend; in the \textit{reflection} stage, the recommender is updated online when its recommendations are rejected, treating the rejected items as negative signal. EAR also moved evaluation to the multi-round setting, in which the system may resume questioning after a rejected recommendation, a protocol that has been standard since. In both CRM and EAR the action space is discrete and enumerates the attribute vocabulary: as \citet{lei2020interactive} note in their characterisation of this family, the policy selects among $|P|+1$ actions — one per candidate attribute plus a recommendation action — which on their evaluation datasets amounts to action spaces of approximately 600 and 8{,}400 dimensions respectively. At that scale, policy learning over the full enumeration already strains sample-efficient training, and this pressure shaped everything that followed.

\subsection{Graph-Based and Hierarchical Successors}

Subsequent work in this lineage concentrated on taming the size of the discrete action space rather than abandoning it, and the dominant tool for doing so has been graph structure over the catalogue. The SCPR framework \citep{lei2020interactive} reframes conversational recommendation as interactive path reasoning on a user--item--attribute graph: confirmed attributes anchor a path through the graph, and candidate next-questions are restricted to attributes adjacent to the current path, which prunes the effective decision set by orders of magnitude. Notably, SCPR's reinforcement learning component is thereby reduced to a \textit{binary} decision — whether to ask or to recommend — while the choice of \textit{which} attribute to ask is delegated to a graph-constrained weighted-entropy heuristic outside the learned policy altogether. This is a revealing design: faced with a large discrete space, the authors found it more effective to learn only the timing decision and to leave question content to a myopic information-theoretic criterion, precisely the separation of concerns that this thesis seeks to overcome by making question content itself learnable.

UNICORN \citep{deng2021unified} re-unified the decision by placing both candidate items and candidate attributes into a single discrete action set scored by a dueling deep Q-network over a dynamically-updated graph representation of the conversation state, refined with graph convolutional encodings of the user, item, and attribute nodes. Tractability is managed through two pruning strategies — preference-based scoring for items and weighted entropy for attributes — that shortlist a small candidate set at each turn from which the Q-network selects. The hierarchical refinements that followed retain the same fundamental structure while subdividing the decision. HutCRS \citep{qian2023hutcrs} organises the interaction over a hierarchical interest tree, first identifying which broad \textit{aspect} the user is interested in before descending to attribute-level questions within that aspect, with a dueling double Q-network scoring decisions over the tree's discrete nodes. DAHCR \citep{zhao2023hierarchical} separates a \textit{director}, which chooses between asking and recommending (trained with a Gumbel-softmax relaxation to keep the discrete choice differentiable), from an \textit{actor}, which selects the concrete attribute or item from the candidate set, with hypergraph-based state representations capturing higher-order relations among the user's confirmed preferences; the two levels provide mutual feedback, with the actor's outcomes informing the director's assessment of its options. CoCHPL \citep{fan2024chain} pushes the hierarchical option framework further still, generating \textit{chains} of choices per turn — sequences of attribute selections produced by intra-option policies under a long-term option policy — explicitly motivated by the observation that ``contemporary CRS are limited to inquiring binary or multi-choice questions based on a single attribute type per round''.

It is worth pausing on what drives this escalating architectural complexity. Each refinement — graph pruning, interest trees, director--actor splits, choice chains — is a response to the same underlying tension: the policy must choose among thousands of discrete possibilities with only a handful of training interactions per dialogue, and so the design burden falls on engineering the candidate set down to a tractable size before the learned component ever sees it. The candidate-generation machinery is heuristic, hand-designed, and catalogue-specific; the learned policy operates only within the shortlist the heuristics permit. An action representation in which proximity encodes semantic relatedness — so that the policy can generalise across related questions rather than treating each as an unrelated arm — is the alternative response to the same tension, and is the route taken in Section~\ref{sec:continuous}.

\subsection{Critical Observations}

Three observations about this lineage are material to the present thesis. First, across all of these systems — from CRM in 2018 to CoCHPL in 2024 — the action space remains discrete, enumerating attributes and items drawn from a structured catalogue; the policy's expressiveness is bounded by that enumeration, and its generalisation across actions is limited by representations learned from the interaction data of a single catalogue. Second, the ``questions'' these systems ask are not natural language in any meaningful sense: they are symbolic attribute queries rendered through templates, evaluated against simulated users that respond symbolically in turn — typically answering attribute queries truthfully from a held-out ground-truth profile and accepting or rejecting recommended items by lookup — with no natural language generation or understanding on either side of the exchange. The linguistic dimension of elicitation, which Chapter~\ref{ch:rl_crs} identified as the principal weakness of pre-LLM CRS, is simply absent from this lineage's experimental designs. Third, the elicitation target is confined to the attribute schema of the catalogue; preferences that users naturally express but that do not correspond to a catalogued facet — tone, mood, plot structure, similarity to other items — are outside the action vocabulary altogether. This matters because studies of natural recommendation dialogues show that human preference talk is dominated by exactly such free-form descriptive language alongside item examples \citep{radlinski2019coached}, and because the qualities users describe most readily (``something dark but funny'', ``a film that takes its time'') are the ones a schema least readily encodes.

The bot-play framework of Chapter~\ref{ch:main} and \citet{makarova2024learning} belongs to this lineage in its training protocol — a learned questioner optimised against the loss reduction of a recommendation model, following the goal-driven self-play paradigm of \citet{das2017learning} — but already extended the action vocabulary to include item attributes alongside items, and demonstrated that the resulting policy concentrates its early questions on the most informative attributes. The natural next step, eliciting preferences over an \textit{open semantic space} of concepts rather than a closed schema, is precisely where the discrete formulation breaks down: the space of concepts worth asking about is unbounded, semantically structured, and shared across catalogues, none of which an enumerated action set can express. This motivates the continuous treatment examined next.

\section{Continuous Action Spaces and Large-Scale Action Grounding}
\label{sec:continuous}

\subsection{Continuous Control and the Deterministic Policy Gradient Family}

Chapter~\ref{ch:rl_crs} reviewed the value-based, policy-gradient, and bandit families of reinforcement learning as they have been applied to CRS; the present section concerns a family largely absent from that literature — continuous control — because it is the foundation of the approach developed in this thesis. Where value-based methods such as deep Q-networks \citep{mnih2015human} require a maximisation over actions at every step, and are therefore confined to action sets small enough to enumerate, the deterministic policy gradient family avoids this maximisation by training an \textit{actor} network that maps states directly to points in a continuous action space, alongside a \textit{critic} that estimates the action-value function and provides the gradient through which the actor improves. Deep Deterministic Policy Gradient (DDPG) \citep{lillicrap2015continuous} established the canonical recipe — replay buffers and target networks borrowed from deep Q-learning, additive exploration noise on the actor's output — and demonstrated continuous control from high-dimensional observations. The family's known weaknesses are equally well documented: DDPG is sensitive to hyperparameters and prone to critic overestimation, failure modes addressed by the twin-critic and delayed-update mechanisms of TD3 \citep{fujimoto2018addressing} and by the maximum-entropy formulation of Soft Actor-Critic \citep{haarnoja2018soft}, which regularises exploration by optimising reward and policy entropy jointly. Two properties of this family bear directly on the present work: its relative sample inefficiency, which the experimental chapters of this thesis confront directly, and the brittleness of deterministic policies in settings where diverse behaviour is itself valuable — a consideration that the entropy-regularised members of the family are designed to address and that the diversity collapse observed in this thesis's early experiments (Chapter~\ref{ch:con}) independently confirms.

\subsection{The Wolpertinger Architecture}

An honest positioning of this thesis requires acknowledging that the core \textit{mechanism} it employs — an actor-critic policy that outputs a continuous vector in an embedding space, which is then grounded to a discrete entity by nearest-neighbour search — is established prior art in reinforcement learning. The Wolpertinger architecture of \citet{dulacarnold2015deep} addresses reinforcement learning in large discrete action spaces by embedding the discrete actions in a continuous space and training a DDPG actor to emit a continuous \textit{proto-action} in that space. The proto-action is then grounded by retrieving its $k$ approximate nearest neighbours among the embedded discrete actions, and — crucially — the critic evaluates each of the $k$ candidates, with the agent executing the candidate of highest estimated value. This propose-then-refine structure serves two purposes. Computationally, it replaces the intractable maximisation over the full action set with a sublinear nearest-neighbour lookup followed by $k$ critic evaluations. Statistically, the critic re-ranking step smooths over the discontinuity introduced by the nearest-neighbour mapping: small movements of the proto-action that would otherwise cause abrupt switches between unrelated actions are mediated by the critic's assessment of the local candidate set. The authors report that small fractions of the action set suffice for $k$, and demonstrate the architecture on tasks with action sets up to one million elements. Importantly for the present thesis, recommender systems were not merely a motivating example: the paper includes a simulated recommendation experiment constructed from a live production recommendation engine, with action sets ranging up to 13{,}138 items, in which the architecture learned effective policies at a fraction of the computational cost of full enumeration.

Two aspects of Wolpertinger deserve emphasis because they inform the design choices examined in this thesis's experimental chapters. First, the action embedding in Wolpertinger is either given a priori or learned from the task; the architecture itself is agnostic to where the embedding geometry comes from, which licenses the substitution made in this thesis of a \textit{pretrained semantic sentence-embedding space} \citep{reimers2019sentence} in which proximity encodes meaning learned from language at large rather than from task interaction data. The practical significance of this substitution is generalisation: a policy operating in a semantic space can, in principle, output proto-actions near concepts it has never been rewarded for, because the geometry already organises related concepts together, and the grounding step can be implemented with standard approximate nearest-neighbour indices \citep{johnson2019billion}. Second, the critic re-ranking step ($k > 1$) is integral to the architecture's stability, not an optional refinement; configurations that ground the proto-action to its single nearest neighbour forfeit the smoothing that makes the critic's learning problem well-posed. This observation proves consequential for the diagnosis, presented in Chapter~\ref{ch:con}, of the training instabilities encountered in this thesis's own continuous-action experiments.

\subsection{Continuous Actions in Recommendation, and What Remains Uncovered}

Within recommender systems research proper, continuous-action formulations have remained rare but are attracting renewed interest. ECoC \citep{wang2024efficient} formulates sequential recommendation as efficient continuous control in a unified low-dimensional action representation abstracted from normalised user and item spaces, arguing that policy learning in the dense representation, followed by grounding to concrete items, is both faster and more transferable than enumeration-based alternatives, and evaluating the approach in offline reinforcement learning settings on interaction logs. The existence of this line of work confirms that the recommender systems community regards continuous-action control as a live methodological option; it also sharpens the boundary of what has actually been attempted.

What this prior art does \textit{not} cover is equally precise. Wolpertinger selects \textit{items to recommend}; it contains no dialogue, no questions, and no preference elicitation. ECoC likewise operates on logged interaction sequences with no conversational component. To the best of the knowledge gathered for this review — including full-text verification of the closest candidates — no published work applies a continuous-embedding action space to the problem of \textit{selecting what to ask} in a conversational recommender; none grounds such actions in a general-purpose semantic sentence-embedding space as opposed to an action embedding learned from the interaction data itself; and none couples the grounded action to a generative language model that verbalises it as a natural question. The contribution of this thesis is therefore not the continuous-action mechanism, which it inherits from \citet{dulacarnold2015deep}, but its application to elicitation question selection over an open vocabulary of semantic concepts, and the integration of that policy with LLM-based question generation — a combination whose absence from the literature was re-verified for this chapter and is further delineated in Section~\ref{sec:gap}.

\section{Large Language Models as Conversational Recommenders}
\label{sec:llm_crs}

\subsection{Zero-Shot Competence}

The publication of \citet{he2023large} marked a turning point for the field. Evaluating GPT-class models as zero-shot conversational recommenders on established CRS benchmarks, the authors found that prompted LLMs, with no task-specific training whatsoever, outperformed all fine-tuned CRS models of the preceding generation. Their accompanying analysis is as important as the headline result. Probing what kind of knowledge the models exploit, they distinguish \textit{content and context knowledge} — facts about items, genres, and the meaning of what the user has said — from \textit{collaborative knowledge}, the patterns of co-consumption that traditional collaborative filtering captures, and find the LLMs' advantage concentrated almost entirely in the former: performance degrades when content cues are removed, and remains comparatively weak on signals that require knowing which items tend to be consumed together. They further document characteristic failure modes of generative recommenders, including a bias toward popular items and a propensity to recommend items outside the evaluation catalogue, which must be controlled by constrained decoding or post-hoc matching. The picture that emerges — fluent, knowledgeable, content-driven recommendation with weak collaborative grounding — directly shapes the hybrid architectures reviewed below, which delegate to traditional recommenders precisely what LLMs lack. It also frames the role taken by LLMs in this thesis: as interface components (question verbalisation, preference extraction, user simulation) wrapped around an explicitly trained recommendation and elicitation core, rather than as the recommendation logic itself.

\subsection{The Evaluation Rethink}

In parallel, \citet{wang2023rethinking} demonstrated that the apparent weakness of LLMs under traditional CRS evaluation protocols was substantially an artefact of those protocols. Established benchmarks score systems on replicating the ground-truth items mentioned at a fixed point in a logged human dialogue; an interactive system that responds to ambiguity by asking a clarifying question, rather than guessing the logged item, is penalised for behaviour that would serve a real user well. Their iEvaLM framework replaces static dialogue replay with interactive evaluation against an LLM-based user simulator that role-plays the user's persona, supporting both free-form conversation and attribute-focused exchanges, and under this protocol the measured performance of LLM-based recommenders improved markedly, in some configurations by large multiples over their replay-based scores. This result simultaneously legitimised LLM-based user simulation as an evaluation instrument — a methodology this thesis adopts in modified form — and exposed how sensitive CRS evaluation is to simulator design: the simulator's knowledge, cooperativeness, and verbosity become experimental parameters that can move results by more than the differences between competing systems. The implications of this sensitivity for measuring elicitation specifically are taken up in Section~\ref{sec:eval_norms}.

\subsection{Architectural Trends: Tool Use, Multi-Agent Decomposition, and Retrieval}

System-building work since 2023 has converged on architectures in which one or more LLMs are orchestrated with external components rather than fine-tuned end-to-end; Chapter~\ref{ch:rl_crs} reviewed the earliest representatives of this pattern (Chat-REC, RecLLM, TALLRec), and the trend has since consolidated along three lines. The first is \textit{multi-agent decomposition}. MACRS \citep{fang2024multi} decomposes the conversational task across cooperating LLM agents — a planner coordinating specialised asking, recommending, and chit-chat agents — and introduces a feedback-aware reflection mechanism operating at two levels: information-level reflection revises the inferred user profile when feedback contradicts it, and strategy-level reflection adjusts the dialogue-act planning itself between turns. MACRec \citep{wang2024macrec} generalises the pattern across recommendation tasks, including conversational recommendation, with manager, reflector, analyst, searcher, and task-interpreter roles collaborating under a deliberate division of labour. The second line is \textit{retrieval augmentation}. CoRAL \citep{wu2024coral} demonstrates that retrieving collaborative evidence — records of which similar users interacted with which items — into the LLM's context substantially improves long-tail recommendation, directly compensating for the collaborative-knowledge weakness identified by \citet{he2023large}; notably, CoRAL learns its retrieval policy by reinforcement learning, treating the composition of the evidence set as a sequential decision problem, an early instance of reinforcement learning returning to the LLM-CRS stack in a component role rather than as the end-to-end controller.

The third line, most relevant to this thesis, reintroduces a \textit{learned decision policy} alongside a frozen LLM. PPDPP \citep{deng2023plug} trains a compact plug-and-play policy language model to select, at each turn, the dialogue strategy that a frozen LLM then realises in natural language; the policy is first supervised on corpora annotated with dialogue acts and then refined by reinforcement learning against LLM-simulated users, with LLM-derived reward signals, demonstrated on negotiation, emotional support, and tutoring dialogues. Most recently, and closest to the present work, RSO \citep{zhao2025reinforced} applies the same decomposition to conversational recommendation itself: a reinforcement-learned planner selects, at each turn, one of thirteen named macro-strategies — including three preference-elicitation strategy types such as opinion inquiry and experience inquiry — and a frozen LLM actor verbalises the chosen strategy, with strategy learning formulated as entropy-regularised policy optimisation against rewards supplied by an LLM judge. RSO's stated motivation is telling: the authors argue that existing LLM-based CRS ``lack explicit optimization of interaction strategies, instead relying on unified prompts and the LLM's internal knowledge to decide how to interact, which can lead to suboptimal outcomes''. The structural similarity to this thesis — learned policy chooses, frozen LLM speaks — is evident; the difference, developed in Section~\ref{sec:gap}, is the granularity and geometry of the action space. RSO selects among a closed inventory of thirteen strategy \textit{types}; the present work selects the semantic \textit{content} of the question from a continuous concept space. The two decisions are complementary rather than competing — a system could in principle use an RSO-style planner to decide \textit{that} elicitation is warranted and a policy of the kind developed here to decide \textit{what} to elicit — but they are not interchangeable, and only the latter addresses the content-selection problem that Sections~\ref{sec:rl_lineage} and~\ref{sec:elicitation_llm} show to be unsolved.

\subsection{Is Strategic Elicitation Solved by LLMs?}

The claim that motivated this thesis — that eliciting preferences strategically is a distinct capability which neither traditional recommenders nor language models possess by default — requires re-examination in the LLM era, and survives it in a sharpened form. The evidence assembled for this review indicates that \textit{prompt-only} LLM recommenders, however fluent, do not reliably implement effective elicitation strategies. \citet{he2023large} found elicitation-related behaviour to be a weak point of zero-shot LLM recommenders relative to their recommendation fluency; the motivation analyses of \citet{zhao2025reinforced} report that unified-prompt systems decide how to interact suboptimally; and user-experience studies of LLM-based CRS report a tendency to recommend prematurely, before sufficient preference information has been gathered \citep{kostric2025should, zhao2025reinforced}. At the same time, this weakness must be stated with care, because it is increasingly specific to \textit{unoptimised} prompting. A growing body of work trains LLMs to ask better questions: \citet{zhang2025modeling} teach LLMs to ask clarifying questions by simulating the future conversation turns that would follow each candidate response, scoring candidates by the eventual outcomes of those simulated futures, and applying direct preference optimisation to the resulting preference labels — in effect, optimising the question through its downstream consequences, the same principle that motivates reward design in this thesis's bot-play framework, realised at the token level rather than in a policy network. RSO, similarly, optimises strategy selection with policy-gradient methods rather than trusting the LLM's defaults.

The honest formulation, adopted henceforth in this thesis, is therefore as follows: strategic preference elicitation remains an open problem for which explicit optimisation — whether in the prompt, the policy, or the action space — measurably outperforms default LLM behaviour; the question of \textit{which} optimisation locus is most effective is unresolved; and, as Section~\ref{sec:eval_norms} argues, it is largely unmeasured, because no published evaluation isolates elicitation strategy from the surrounding system. This formulation directly motivates the evaluation contribution proposed in Chapter~\ref{ch:con}.

\section{Preference Elicitation Beyond Recommendation Dialogue}
\label{sec:elicitation_llm}

Preference elicitation has a research history older than CRS, rooted in decision theory, where the object of inquiry is a utility function and questions are chosen to reduce uncertainty about it — whether by minimax regret over the remaining feasible utilities \citep{viappiani2009regret}, by Bayesian updating of a belief over user models \citep{mangili2020bayesian}, or by bandit-style absolute and comparative questions about items \citep{christakopoulou2016towards}. The LLM era has revived this tradition in forms directly relevant to this thesis, because the language model removes what was previously the binding constraint: the need to confine questions to forms a system could both generate and interpret.

GATE \citep{li2023eliciting} demonstrates that LLMs can themselves conduct free-form preference interviews. In their framework, the language model elicits a user's intended task specification through interaction — generating open-ended questions, probing edge cases, and proposing example inputs for the user to judge — across domains that include content recommendation, and the elicited specifications improve downstream decisions over both user-written prompts and labelled examples, while often requiring less user effort. GATE is, in effect, an existence proof that LLMs can \textit{conduct} elicitation; what it deliberately does not provide is a learned strategy — its questioning is generated greedily by prompting, with no optimisation target tied to a downstream recommender and no notion of multi-turn planning. PEBOL \citep{austin2024bayesian} supplies the principled decision-theoretic treatment that GATE lacks, framing natural-language preference elicitation as Bayesian optimisation: the system maintains a probabilistic belief over the utilities of catalogue items, and at each turn an acquisition strategy in the style of Thompson sampling or upper-confidence selection identifies the item description about which a yes/no preference query would be most valuable, with the LLM serving to verbalise the query and to map the user's natural-language answer back into the belief update via entailment-style inference. PEBOL outperforms monolithic prompted-GPT elicitation, which is itself notable as a direct, controlled demonstration that adding explicit decision-theoretic structure on top of an LLM beats trusting the LLM's own conversational instincts. Its limitations mirror its strengths: the acquisition is deliberately myopic — a one-step criterion rather than a learned multi-turn policy — and it operates over item-level descriptions rather than an open concept space.

Within the CRS mainline, the elicitation question itself has come under study as an object of design. \citet{kostric2024generating} generate \textit{usage-oriented} elicitation questions — asking how or in what context an item will be used, rather than which attributes are desired — on the grounds that users with limited domain expertise can answer usage questions more reliably than schema questions; their generation pipeline mines item reviews for usage signals, and the work reflects a broader shift toward elicitation grounded in user context rather than catalogue schemata. Complementing the question-content perspective, \citet{kostric2025should} examine the conversational \textit{style} of elicitation, contrasting a fast-paced, direct, high-involvement style with a patient, accommodating, high-considerateness style in a controlled user study, and find that style measurably affects satisfaction and elicitation effectiveness, with the ability to adapt style to the user performing best. Both results underline that elicitation quality is multi-dimensional; both also concern dimensions orthogonal to the one this thesis optimises, namely the strategic selection of what to ask next.

Taken together, this strand confirms both halves of the thesis premise: elicitation quality measurably matters and is actively studied, yet none of these works \textit{learns a multi-turn questioning strategy} — GATE and usage-question generation are unoptimised generation, PEBOL is single-step Bayesian acquisition, and the style work concerns the form rather than the content of questions. The learned-strategy tradition surveyed in Section~\ref{sec:rl_lineage} and the LLM-era elicitation work surveyed here have, as of this writing, intersected only at the coarse strategy-type level \citep{zhao2025reinforced, deng2023plug}; the content-level intersection remains open.

\section{Evaluation Methodology in the LLM Era}
\label{sec:eval_norms}

\subsection{User Simulation}

Evaluation practice has shifted alongside system architecture. The discrete-action lineage of Section~\ref{sec:rl_lineage} evaluated against rule-based simulators that answer attribute queries truthfully from a held-out profile and accept a recommendation exactly when it matches ground truth \citep{lei2020estimation, deng2021unified} — fully verifiable, but linguistically empty and behaviourally rigid. The use of LLMs as user simulators, introduced systematically by \citet{wang2023rethinking}, has since become the dominant protocol for interactive evaluation, offering natural language, persona conditioning, and varied behaviour. The risks, acknowledged but not yet resolved in the literature, are equally real: the simulator may diverge from human behaviour in unknown ways; its underlying model carries biases (toward popular items, toward agreeableness) that interact with the system under test; its knowledge of the full catalogue can leak information a real user would not provide; and a circularity arises when the same model family simulates the user, generates the system's utterances, and judges the outcome. Simulator design has thus become part of the experimental design, and results are only as interpretable as the simulator is controlled.

This thesis's evaluation design — simulated users instantiated from held-out MovieLens profiles \citep{harper2015movielens}, with the simulator constrained to answer from the profile's ground-truth ratings, optionally through tool-mediated lookups that make every simulated statement auditable against the profile — sits deliberately between the rule-based and free-LLM extremes. It trades some linguistic naturalism for verifiable consistency between simulated statements and ground truth, which is essential when the measured quantity is information gain: an unconstrained simulator that hallucinates a preference injects unmeasurable noise into precisely the signal under study.

\subsection{Measuring Elicitation, and a Methodological Gap}

A further observation from this review is methodological. While recommendation accuracy metrics are thoroughly standardised — NDCG, MRR, hit rate at fixed cut-offs, success rate and average turns for multi-round dialogues — there is no established benchmark or agreed protocol for measuring \textit{strategic elicitation ability} in isolation: that is, how much recommendation-relevant information a system's questions extract per turn, holding the recommender, the simulator, and the candidate pool fixed. Published comparisons embed elicitation inside whole-system evaluations, where its effect is confounded with the quality of recommendation, generation, and language understanding; a system may score well because it asks good questions, or because its recommender squeezes more from the same answers, and the published metrics cannot tell these apart. The confound matters practically as well as conceptually, because the candidate policies now in play are of radically different kinds — graph-pruned Q-networks, Bayesian acquisition, prompted LLMs, policy-plus-verbaliser hybrids — and the field has no instrument for comparing them on equal terms.

The diagnostic experiments conducted in this thesis (Chapter~\ref{ch:con}) add a quantitative warning to this picture: commonly used reward signals for elicitation, such as per-turn changes in ranking metrics over modest candidate pools, can be statistically too noisy to support either learning or fair comparison — in this thesis's measurements, per-turn metric deltas were an order of magnitude smaller than their per-user standard deviation under realistic pool sizes and evaluation filters. Any evaluation (or training) protocol for elicitation must therefore be designed around the signal-to-noise ratio of its measurement instrument, a consideration absent from the published evaluations reviewed here. The absence of such a protocol is itself a gap this thesis is positioned to address, and grounds the evaluation-first contribution proposed in Chapter~\ref{ch:con}.

\section{The Updated Research Gap}
\label{sec:gap}

Table~\ref{tab:positioning} summarises the positioning of the works most relevant to this thesis along the dimensions established in Section~\ref{sec:problem}: whether a questioning strategy is learned, the structure of the action space, how questions are realised in language, and the role of LLMs.

\begin{table}[ht]
\centering
\caption{Positioning of closest related work. ``Strategy learned'' refers to multi-turn optimisation of what to ask; ``action space'' refers to the elicitation decision.}
\label{tab:positioning}
\small
\begin{tabular}{p{3.4cm}p{2.5cm}p{3.2cm}p{2.6cm}p{2.2cm}}
\toprule
\textbf{Work} & \textbf{Strategy learned} & \textbf{Action space} & \textbf{Question realisation} & \textbf{LLM role} \\
\midrule
CRM, EAR \citep{sun2018conversational, lei2020estimation} & RL & Discrete ($|P|{+}1$ attributes) & Template & None \\
SCPR \citep{lei2020interactive} & RL (binary) & Ask/recommend; attribute via entropy heuristic & Template & None \\
UNICORN; HutCRS; DAHCR; CoCHPL \citep{deng2021unified, qian2023hutcrs, zhao2023hierarchical, fan2024chain} & RL & Discrete pruned/hier-archical candidate sets & Template/symbolic & None \\
Bot-play elicitation \citep{makarova2024learning} & RL (REINFORCE) & Discrete (items + attributes) & Symbolic vocabulary & None \\
Wolpertinger \citep{dulacarnold2015deep}; ECoC \citep{wang2024efficient} & RL & \textbf{Continuous} embedding $\to$ kNN & n/a (item selection, no dialogue) & None \\
GATE \citep{li2023eliciting}; usage questions \citep{kostric2024generating} & None (prompted/ generated) & Open language & LLM-generated & Generator \\
PEBOL \citep{austin2024bayesian} & Bayesian, myopic & Item-utility queries & LLM-generated & Acquisition fn. \\
PPDPP \citep{deng2023plug}; RSO \citep{zhao2025reinforced} & RL & Discrete strategy types (13 in RSO) & Frozen LLM verbalises & Verbaliser, judge, simulator \\
\citet{zhang2025modeling} & DPO (token-level) & Token sequences & LLM-generated & Policy itself \\
\midrule
\textbf{This thesis (proposed)} & RL (actor-critic) & \textbf{Continuous semantic embedding $\to$ entity (kNN)} & \textbf{LLM verbalises grounded concept} & Verbaliser, extractor, simulator \\
\bottomrule
\end{tabular}
\end{table}

The gap claimed in earlier versions of this work — that no existing system combines reinforcement learning in continuous action spaces, LLM-based question generation, and preference elicitation — survives the verification conducted for this chapter, but only when stated as a claim about the \textit{combination}, with each constituent honestly attributed. Continuous-action reinforcement learning with nearest-neighbour grounding is the Wolpertinger architecture \citep{dulacarnold2015deep}, demonstrated on recommendation-scale action sets; continuous control has been applied to sequential recommendation \citep{wang2024efficient}; the policy-chooses/LLM-speaks decomposition exists at the strategy-type level \citep{deng2023plug, zhao2025reinforced}; and LLMs both generate elicitation questions \citep{li2023eliciting, kostric2024generating} and learn, at the token level, to ask clarifying questions \citep{zhang2025modeling}. What no verified work does is learn a multi-turn elicitation policy whose actions are points in a continuous \textit{semantic} space — decoupling the policy's decision (``what concept to probe next'') from both the catalogue schema that constrains the discrete lineage and the closed strategy inventories of RSO and PPDPP — grounded to interpretable entities and verbalised by an LLM. Equally, no verified work provides a controlled protocol for measuring strategic elicitation ability across approaches of such different kinds; the lineage systems, the prompt-based LLM systems, and the Bayesian systems have never been compared on equal terms with the recommender, simulator, and evaluation pool held fixed.

These two gaps jointly define the research questions for the remainder of this thesis:

\begin{itemize}
    \item \textbf{RQ1 (measurement).} How can strategic preference elicitation be measured in isolation — with the recommendation model, user simulator, and candidate pool held fixed — such that policies of heterogeneous kinds (random, information-theoretic, prompted-LLM, and learned) can be compared on equal terms, and what signal-to-noise properties must the measurement instrument satisfy for such comparisons to be statistically meaningful?
    \item \textbf{RQ2 (the LLM baseline).} Under such a protocol, how strategically do prompt-only LLM questioners actually behave relative to principled non-learned baselines such as greedy information gain — quantifying, rather than asserting, the elicitation weakness attributed to LLMs in the literature?
    \item \textbf{RQ3 (the learned policy).} Does a multi-turn policy learned in a continuous semantic action space, grounded to entities and verbalised by an LLM, outperform both the discrete formulation of \citet{makarova2024learning} and the strongest baselines from RQ2 — and under what training regime (sample budget, reward signal, grounding strategy) does the continuous formulation become learnable at all?
\end{itemize}

This thesis accordingly advances two contributions in its remaining chapters. The first is methodological, answering RQ1 and RQ2: a controlled evaluation framework for strategic preference elicitation, in which a fixed, polarity-aware recommendation model serves as the measurement instrument, simulated users answer verifiably from held-out ground truth, and elicitation policies of heterogeneous types — random and popularity baselines, greedy information-gain, prompt-only LLM questioning, and learned policies including the published bot-play approach \citep{makarova2024learning} — are compared on information gained per conversational turn. The second is algorithmic, answering RQ3: the extension of the bot-play framework to a continuous semantic action space with LLM verbalisation, evaluated within that framework against the strongest non-learned and LLM-based baselines. The diagnostic findings already obtained (Chapter~\ref{ch:con}) — concerning reward signal-to-noise ratios, the necessity of separating preference polarity from semantic content in user-state representations, and the sample-efficiency demands of continuous-action learning — directly inform the design of both contributions.
